\section{Objetivos}

\begin{itemize}
	\item \textbf{Comprender y aplicar} los principios fundamentales de la resonancia eléctrica en circuitos de corriente alterna para identificar las condiciones de operación en las que las reactancias se anulan mutuamente.
	
	\item \textbf{Determinar} analíticamente, mediante el uso de modelos matemáticos y leyes fundamentales de circuitos, la frecuencia de resonancia ($\omega_0$), las frecuencias de media potencia ($\omega_1, \omega_2$) y el factor de calidad ($Q$) del sistema.
	
	\item \textbf{Validar}, a través de herramientas de simulación y el análisis de los videos de la sesión en línea, el comportamiento de la corriente y la impedancia en el circuito serie propuesto.
	
	\item \textbf{Verificar} los niveles de potencia disipada en las frecuencias características ($\omega_0$, $\omega_1$ y $\omega_2$) para contrastar los resultados teóricos con la evidencia experimental obtenida en el laboratorio.
	
	\item \textbf{Analizar} el efecto de la selectividad del circuito sobre la respuesta del sistema, centrando la atención en los detalles esenciales para la correcta aplicación de los principios fundamentales de la ingeniería eléctrica.
\end{itemize}